\subsection{Изменение настроек}

К функционалу изменения настроек программного средства можно отнести следующие компоненты:

\begin{itemize}
    \item получение настроек;
    \item изменение настроек;
\end{itemize}

\subsubsection{Получение настроек}

Получение настроек начинается с открытия ключа реестра Windows, в котором хранятся сохранённые параметры конфигурации приложения.

После успешного открытия ключа реестра функция последовательно читает три параметра: тему оформления, размер шрифта и имя шрифта. Для каждого параметра выполняется проверка успешности чтения и валидация значения.

При чтении темы оформления проверяется, что значение находится в допустимом диапазоне (от System до Dark), и если значение валидно, оно устанавливается в глобальную переменную.

При чтении размера шрифта проверяется, что значение находится в диапазоне от 8 до 72 пунктов, и если значение валидно, оно устанавливается в глобальную переменную.

При чтении имени шрифта проверяется успешность чтения строкового значения, и если чтение успешно, имя шрифта устанавливается в глобальную переменную.

После обработки всех параметров ключ реестра закрывается, и вызывается функция UpdateThemeColors для применения загруженных настроек темы к цветовой схеме интерфейса.

Если в процессе работы алгоритма произошла ошибка (например, ключ реестра не существует, доступ запрещён или значение невалидно), функция использует значения по умолчанию (тема System, размер шрифта 12, имя шрифта "Segoe UI") и продолжает работу.

\begin{lstlisting}[style=CodeListing]
void Load() {
    HKEY hKey;
    if (RegOpenKeyExW(HKEY_CURRENT_USER, L"Software\\KursWork\\FileManager", 
                      0, KEY_QUERY_VALUE, &hKey) == ERROR_SUCCESS) {
        DWORD theme = (DWORD)ThemeChoice::System; 
        DWORD sz = sizeof(theme);
        if (RegGetValueW(hKey, nullptr, L"Theme", RRF_RT_REG_DWORD, 
                        nullptr, &theme, &sz) == ERROR_SUCCESS) {
            if (theme <= (DWORD)ThemeChoice::Dark) {
                g_themeChoice = (ThemeChoice)theme;
            }
        }
        DWORD size = 12; 
        sz = sizeof(size);
        if (RegGetValueW(hKey, nullptr, L"FontSize", RRF_RT_REG_DWORD, 
                        nullptr, &size, &sz) == ERROR_SUCCESS) {
            if (size >= 8 && size <= 72) {
                g_fontPt = (int)size;
            }
        }
        wchar_t buf[128]; 
        DWORD bytes = sizeof(buf);
        if (RegGetValueW(hKey, nullptr, L"FontFace", RRF_RT_REG_SZ, 
                        nullptr, buf, &bytes) == ERROR_SUCCESS) {
            g_fontFace = buf;
        }
        RegCloseKey(hKey);
    }
    UpdateThemeColors();
}
\end{lstlisting}

\subsubsection{Изменение настроек}

Изменение настроек начинается с обработки команды пользователя из меню приложения. В зависимости от выбранной команды определяется тип изменяемой настройки (тема, имя шрифта или размер шрифта).

При изменении темы оформления проверяется идентификатор выбранной команды меню, и в зависимости от него устанавливается соответствующее значение темы (System, Light или Dark). После установки темы вызывается функция Save для сохранения изменений в реестр, и пользователю отображается сообщение о необходимости перезагрузки приложения.

При изменении имени шрифта проверяется идентификатор выбранной команды меню, и в зависимости от него устанавливается соответствующее имя шрифта ("Segoe UI", "Consolas" или "Courier New"). После установки имени шрифта вызывается функция ApplyFontToControls для применения нового шрифта ко всем элементам интерфейса, обновляются отметки в меню, и вызывается функция Save для сохранения изменений.

При изменении размера шрифта проверяется идентификатор выбранной команды меню или клавиатурной комбинации. Если это команда увеличения размера, проверяется, что новый размер не превышает 72 пункта, и размер увеличивается на 1. Если это команда уменьшения размера, проверяется, что новый размер не меньше 8 пунктов, и размер уменьшается на 1. После установки размера шрифта вызывается функция ApplyFontToControls для применения нового размера ко всем элементам интерфейса, обновляются отметки в меню, и вызывается функция Save для сохранения изменений.

\begin{lstlisting}[style=CodeListing]
case ID_MENU_THEME_SYSTEM:
case ID_MENU_THEME_LIGHT:
case ID_MENU_THEME_DARK: {
    if (LOWORD(wParam) == ID_MENU_THEME_SYSTEM) {
        Settings::SetTheme(ThemeChoice::System);
    } else if (LOWORD(wParam) == ID_MENU_THEME_LIGHT) {
        Settings::SetTheme(ThemeChoice::Light);
    } else {
        Settings::SetTheme(ThemeChoice::Dark);
    }
    Settings::Save();
    return 0;
}
case ID_MENU_FONT_SEGOE:
case ID_MENU_FONT_CONSOLAS:
case ID_MENU_FONT_COURIER: {
    if (LOWORD(wParam) == ID_MENU_FONT_SEGOE) {
        Settings::SetFontFace(L"Segoe UI");
    } else if (LOWORD(wParam) == ID_MENU_FONT_CONSOLAS) {
        Settings::SetFontFace(L"Consolas");
    } else {
        Settings::SetFontFace(L"Courier New");
    }
    UI::ApplyFontToControls(hwnd, *controls);
    UI::UpdateMenuChecks(hwnd);
    Settings::Save();
    return 0;
}
case ID_CMD_FONTSIZE_INC: {
    Settings::SetFontSize(std::min(72, Settings::GetFontSize() + 1));
    UI::ApplyFontToControls(hwnd, *controls);
    UI::UpdateMenuChecks(hwnd);
    Settings::Save();
    return 0;
}
case ID_CMD_FONTSIZE_DEC: {
    Settings::SetFontSize(std::max(8, Settings::GetFontSize() - 1));
    UI::ApplyFontToControls(hwnd, *controls);
    UI::UpdateMenuChecks(hwnd);
    Settings::Save();
    return 0;
}
\end{lstlisting}

Если в процессе работы алгоритма произошла ошибка (например, не удалось сохранить настройки в реестр), изменения применяются только к текущей сессии приложения и не сохраняются для следующего запуска.